\section{Introduction}

An odd cycle in the permutation induced by two lines of a Latin square
supports a particularly simple Latin trade: swap the two lines on that cycle.
The move preserves the Latin property.  In the row and column views it
reverses the conventional row--column Latin sign; in the symbol view it
preserves that sign but reverses the two parastrophically transported signs.
It is therefore a natural primitive for attempts to construct local
sign-reversing involutions in parastrophic versions of the Alon--Tarsi parity
setting.  B\'erczi's
recent experimental program explicitly studies such local corrections as
candidate ingredients for global constructions~\cite{Berczi2026}.  That work
is exploratory and formulates conjectural mechanisms rather than claiming
proofs.

This paper isolates the underlying existence question.  For each of the row,
column, and symbol views, we ask whether some pair of lines induces a
permutation with a nontrivial odd cycle.  Recording the three answers gives a
pattern in $\{F,T\}^3$.  A square of pattern $\FFF$ has no such local
odd-cycle move in any view; equivalently, it avoids this trade primitive for
all three pairwise parity products.  This distinction matters: a symbol-view
odd trade need not reverse the fixed conventional row--column sign.  Cycle
structures of Latin-square lines have long been useful for parity and
symmetry questions; see, for example, Kotlar~\cite{Kotlar2012}.

Our first contribution is structural.  Failure patterns are componentwise
monotone from Latin subsquares to their ambient squares.  They are also
exactly preserved by a Latin quotient when all congruence blocks have size
two, while any nontrivial odd congruence-block size forces pattern TTT.  It
follows that a hypothetical FFF quasigroup of order $2p$, with $p$ an odd
prime, must be congruence-simple in every isotope.  Iteration proves the
power-of-two conjecture for quasigroups admitting a prime-factor congruence
series.  A twisted binary extension of $C_4$ also gives an explicit,
census-free non-group-isotopic FFF loop of order $8$.  This excludes every
congruence-imprimitive order-$10$ quasigroup construction.  A later
one-view reduction separately excludes imprimitive generated line actions;
the exhaustive palette masters developed below subsequently exclude the
remaining FFF case altogether.

A second subsquare theorem identifies the remaining binary line-incidence
defect exactly.  Nonuniversal vectors in the left binary line kernel are in
bijection with half-order subsquares, and quotienting by the two universal
constant relations groups them into complementary four-packs.  This gives a
closed counting formula and forces maximal binary line rank for every FFF
order congruent to $2$ modulo $4$ and greater than $2$.  At order $10$ this
is a sharp necessary condition, but not a decision: the tracked non-FFF
diagnostic tables already have the same maximal rank.

Specializing Moorhouse's $3$-net character formula, every additional
prime-field line dependency
defines a balanced three-sorted cyclic quotient homotopy; after isotopy its
fibres form a quasigroup congruence.  It follows that every Smith torsion
prime divides the square order.  Combining this with the congruence-pattern
theorem shows, without a primitive-group classification, that one F view at
order $2p$ for odd prime $p$ forces the entire integer line-incidence lattice
to be saturated.  Thus a hypothetical order-$10$ FFF square has rank $28$
over every field, although current non-FFF diagnostics already satisfy this
stronger necessary condition.

We additionally derive a cycle-parity condition satisfied by every FFF line
pair.  It is equivalent to the even-cycle condition at degree $6$ and leaves
only one odd-cycle type, $5+3+2$, at degree $10$.  An exact order-$8$ scan
shows that the condition is genuinely weaker than FFF, but it provides a
smaller target for certified order-$10$ obstruction searches.

The exact even-cycle condition itself admits a second formulation.  A
permutation has only even cycles exactly when it preserves no odd-cardinality
subset.  Summing the induced $k$-subset actions of a Latin line family
therefore gives a zero-one matrix for every relevant odd $k$ exactly when the
view is F.  This yields both a structural lift characterization and a direct
proof-producing encoding; at order $10$, only the levels $k=3$ and $k=5$ are
needed.

The local alternating bipartitions also glue globally.  On the dual of a
canonical three-coloured flag surface, assigning one basis vector to every
edge colour defines an $\mathbb F_2^3$-valued cochain whose colour-stratified
coboundary detects the complete view pattern.  It is a cocycle exactly for FFF
squares; then its three coordinate classes sum to the first
Stiefel--Whitney class and force equal holonomy fibres and divisibility of
every component's flag count.  This gives a topological realization of
the slice-wise even-cycle condition rather than another scalar aggregate.

For arbitrary Latin squares $L$ and $K$ we further prove
\[
  \pat(L\times K)=\pat(L)\mathbin{\mathrm{OR}}\pat(K)
\]
componentwise.  We also prove
\[
  L\times K\text{ is group-isotopic}
  \quad\Longleftrightarrow\quad
  L\text{ and }K\text{ are both group-isotopic}.
\]
The proofs are short but useful: the induced permutations are Cartesian
products, and their cycle lengths are least common multiples of factor cycle
lengths.  The group-isotopy statement follows from a basis-family closure
criterion.

Our second contribution is an exhaustive order-$8$ computation.  McKay,
Meynert, and Myrvold enumerated isotopy and main classes of small Latin
squares~\cite{McKayMeynertMyrvold2007}; the associated ANU data page supplies
one representative of each of the $283657$ order-$8$ main
classes~\cite{McKayLatinData}.  Scanning those representatives gives the full
three-view pattern distribution and exactly $230$ FFF main classes.  Only
$5$ are group-isotopic; the remaining $225$ are not.  The frozen input used
for our computation was downloaded again from the cited ANU URL and matched
byte for byte, with SHA-256
\begin{center}
\small\nolinkurl{e5980ff5bc4ac46cbbefd1f1333258d0905724c42a976c4afe61ea501a4ccf87}.
\end{center}

Combining the structural and computational parts yields infinite families.
Every order $2^m$, $m\geq 3$, realizes all eight patterns, and every such
order contains a non-group-isotopic FFF square.  The latter existence claim
now follows entirely theoretically from the explicit twisted order-$8$ loop;
only the realization of all eight patterns uses the census.  We make no claim
that the product construction is injective on isotopy classes or main classes.

Finally, exhaustive generation of all $9408$ reduced squares of order $6$
finds no FFF square.  The generated set agrees exactly with the $9408$-square
ANU reduced-order-$6$ file.  At order $10$, two exhaustive cycle-palette
master searches prove nonexistence.  Together with the elementary odd-order
obstruction, these results verify through order $10$ the conjecture that FFF
squares of order greater than one exist exactly at powers of two.  The general
conjecture remains open; its first undecided even non-power-of-two order is
$12$.

The paper does not claim a proof of the Alon--Tarsi conjecture.  It establishes
the precise limitations of one local-trade primitive, gives exact finite
data, and provides a product construction that converts the order-$8$
phenomenon into infinite families.

The exposition places the structural theory, products, and finite census
before the denser order-$10$ diagnostic machinery.  The flag-surface and
pair-label constructions are retained as a second research layer rather than
as prerequisites for the core FFF results.
